\documentclass[10pt,a4paper]{amsart}
\usepackage[margin=19mm]{geometry}
\usepackage{amsmath,amssymb,mathtools}
\usepackage[scr=rsfs,cal=euler]{mathalfa}
\usepackage{xfrac}
\usepackage[unicode,hidelinks,bookmarks=false]{hyperref}
\newcommand{\ie}{i.e.\ }
\newcommand{\cf}{cf.\ }
\newcommand{\resp}{respectively\ }
\newcommand{\Subsection}{\subsection}
\providecommand{\nobreakdash}{\nobreak\hbox{-}}
\setlength{\emergencystretch}{2em}
\multlinegap0pt
\usepackage{amssymb}
\usepackage{dsfont}
\usepackage{bbm}
\usepackage{xcolor}
\newtheorem{lemma}{Lemma}
\newtheorem{proposition}{Proposition}
\newtheorem{theorem}{Theorem}
\newcommand{\Uc}{\mathbf{UC}}
\newcommand\bR{{\mathbb{R}}}
\newcommand\ra{\rightarrow}
\title{Partially hyperbolic flows on flat vector bundles with an application to complete affine manifolds.}
\author{Suhyoung Choi}
\date{Source: arXiv:2509.26117v2 (CC BY 4.0)}
\begin{document}
\maketitle

\noindent\emph{Source and licence.} Partially hyperbolic flows on flat vector bundles with an application to complete affine manifolds.. Version arXiv:2509.26117v2 (CC BY 4.0), distributed by arXiv under the Creative Commons Attribution 4.0 International licence, http://creativecommons.org/licenses/by/4.0/, as recorded in the arXiv licence metadata for this exact version and shown on the abstract page https://arxiv.org/abs/2509.26117v2. Full licence notice: \texttt{licenses/arxiv-2509.26117-CC-BY-4.0.txt}.\par\smallskip

\noindent\textit{Editorial scope.} Counted unit: the article's theorem thm:identify (source0.tex lines 2200--2208). The article's complete original proof of it is delivered with it (source0.tex lines 2209--2223, one \texttt{proof} environment of 15 lines). No mathematics is written, repaired or re-derived: the statement and the proof are the article's own lines, in the article's own notation, and no retained line is substituted or re-flowed.  Retained uncounted as the source-local context the proof needs: lines 1914--1919; lines 2072--2083.  Nothing was omitted from any retained span, so the package carries no omission record.  No retained line contains a citation, so the wrapper carries no bibliography and no bibliography entry.  No retained line contains a figure environment, an image-inclusion command or drawing code, so no image is delivered and the package declares no asset dependency.  Editorial formatting only, in addition: an AMS \texttt{amsart} preamble that restates the article's own declarations and macros used by the retained lines, namely the environments \texttt{lemma}, \texttt{proposition}, \texttt{theorem} on separate counters, together with the standard packages those lines need.  The retained environments print in this document as Lemma 1, Proposition 1, Theorem 1 in order of appearance, whereas the article prints the same items with the numbering of its own full document.  The complete native source of the article as distributed in the arXiv e-print for this exact version is bundled unchanged as \texttt{sources/186058/source0.tex}.
\begin{lemma} \label{lem:identify} 
	Suppose that $\Gamma_M$ is Gromov hyperbolic. 
	Then there is a quasi-isometric homeomorphism $\mathcal{F}:\Uc \hat M \ra G\hat M$
	by taking a unit vector $\vec{u}$ at $x$ to a complete isometric geodesic 
	$\bR \ra \hat M$ passing $x$ tangent to $\vec{u}$. 
	\end{lemma} 

\begin{proposition} \label{prop:Uc} 
	Let $M$ be a compact manifold with a covering map $\hat M \ra M$
	with a deck transformation group $\Gamma_M$. 
	Suppose that $\Gamma_M$ is word-hyperbolic.

	Then there is 
	a quasi-isometric surjective map 
	$\mathcal{E}: G\hat M \ra \partial^{(2)}_\infty \Gamma_M \times \bR$ with compact fibers. 
	%Here $C = 4\delta$ if $\hat M$ is $\delta$-hyperbolic, and 
	The map obtained from composing with 
	a projection to the first factor is a map given by taking the endpoints of complete isometric geodesics. 
\end{proposition} 

\begin{theorem}\label{thm:identify} 
	%	There is a metric on $\partial^{(2)} M \times \bR$ so that 
	The map	
	$\mathcal{E}\circ \mathcal{F}:\Uc \hat M \ra \partial_\infty^{(2)} \Gamma_M \times \bR$ 
	is a quasi-isometry where each complete isometric 
	geodesic in $\Uc \hat M$ goes to 
	$(t_1,t_2)\times \bR$ for its pair $(t_1, t_2)\in \partial_\infty^{(2)} \Gamma_M$ 
    of endpoints. 
\end{theorem} 

\begin{proof} 
	Proposition \ref{prop:Uc} and Lemma \ref{lem:identify} imply
	the result. %\qed
	% 	We construct a metric on  $\partial^{(2)} M \times \bR$ the space 
	% 	as follows: 
	% 	For each pair of complete isometric geodesics going to the same 
	% 	sets of form $(t_1, t_2)\times \bR$, we use one of 
	% 	the metric restricted to it. 
	% 	For pairs of points $\partial^{(2)} M \times \bR$
	% 	in different factors of $\bR$, 
	% 	we take the inverse images of $\mathcal{E}\circ \mathcal{F}$ of 
	% 	the pair of points and taking the maximal distance in $\Uc \hat M$. 
	% 	The map restrict to a complete isometric geodesic to an $\bR$-factor is an
	% 	isometry by construction in the proof of Proposition \ref{prop:Uc}. 
\end{proof} 

\end{document}
